\chapter{How the literature was searched}
\label{ch:lit-search}

% Author's note (verification and reproduction):
%   sessions      : screenlit-q1 .. screenlit-q4 (one per question) and
%                   screenlit-q1-top .. screenlit-q4-top (the first 100
%                   places of the OpenAlex and arXiv answers), private
%                   data directory of the review of 2026-11-23
%   queries       : lit-screening/work/queries.tsv, queries-added.tsv
%   provider check: lit-screening/work/providers-check-2.txt
%   checks        : lit-screening/work/checks/q*.json (sessions check)
%   decisions     : lit-screening/work/decisions/*.tsv, dec_q*.py
%   samples       : lit-screening/work/tier2-sample-q*.tsv (seed
%                   20261123+q), sample60-q2/q3.txt (seed 20261123)
%   report        : lit-screening/work/report.md
%   tables        : doc/methods/tables/*.tex, made from copies of the
%                   sessions with this branch's scholar (matrix cells as
%                   hits_label gives them; audit tables by
%                   generate_audit_table with --bearing-only and
%                   --cited-from doc/bibliography.bib)

The four chapters that follow each answer one question with the
literature.
They share one search, made on 2026-11-23 with Scholar by one agent,
which this chapter describes once; each chapter then gives its own
queries, their hits per database, its tiers and its sample, and the
records it judged to bear on its question.

\section{Questions and claims}
\label{sec:search-questions}

Each question was put as the claim that the procedure would rest on if
the claim held, so that the searches could find evidence against it as
well as for it.

\begin{description}
\item[First question (\cref{ch:stopping})] Stopping a prioritised screen
  before the end of the list is reliable enough: recall targets and
  stopping rules keep the share of missed relevant studies small.
\item[Second question (\cref{ch:abbreviated})] An abbreviated search or
  screen (fewer databases, a result limit, a rapid review, single
  screening) loses few included studies and leaves pooled estimates and
  conclusions unchanged.
\item[Third question (\cref{ch:databases})] A few bibliographic databases
  (Scopus, Web of Science, Google Scholar, OpenAlex, Semantic Scholar), or
  the first results of a ranked source, find nearly all the relevant
  records a review needs.
\item[Fourth question (\cref{ch:llm-screening})] A cheaper first pass
  with a language model (title-only, a small model, a model on the
  confident band) misses few relevant studies compared with human
  screening.
\end{description}

\section{Databases and queries}
\label{sec:search-databases}

\texttt{scholar providers check}, run before the searches, found
OpenAlex, Scopus (with the full record view), Web of Science and IEEE
Xplore working, Semantic Scholar and arXiv answering with rate limits
(HTTP 429) at that moment, and DBLP failing: it answered with an HTML page
instead of data.
DBLP was therefore not searched.
Semantic Scholar stayed rate-limited through the day, and its answers to
four queries of the third question are incomplete; arXiv answered every
search.

Each question has one session, and every query of a question ran on the
same six databases with the same keywords and without a limit:
OpenAlex in title and abstract (\texttt{-{}-in title-abstract}), Semantic
Scholar through its bulk search, Scopus as \texttt{TITLE-ABS-KEY(\dots)},
Web of Science as \texttt{TS=(\dots)}, IEEE Xplore and arXiv.
The sessions record all of this per search (no limit, the title and
abstract filter for OpenAlex, the bulk search for Semantic Scholar), and
every search is dated 2026-11-23.
The query tables of the following chapters give the queries verbatim and
the hits per database as \texttt{scholar sessions show} reports them.
A zero from IEEE Xplore or arXiv is an answer, not a failure: both
answered other queries of the same syntax with hits.
PubMed is not among Scholar's databases; the records from health methods
journals, where most of the evidence turned out to be, came through
OpenAlex, Scopus and Web of Science (the database column of the tables
of bearing records shows which).

Each question has support queries, at least one second support wording,
and counter queries that reuse every concept term of the support
queries with words for failure, missed studies or changed conclusions.
Some queries were added after the check by meaning
(\cref{sec:search-check}); the query tables mark them.
In all, the four questions have 35 queries.

Before the searches, a few studies known to bear on each question were
named, by following the references of two reviews (one on stopping
rules in education, one on database combinations) through OpenAlex.
All of them came back except two: one study on the contribution of
databases was returned only by the second question's queries, though it
bears on the third too, and one older study on searching several
databases was returned by no query and is not cited.

Scholar deduplicated each search's records: records with one DOI, or
without one but with the same title and first author and years at most
one apart, were merged, and a preprint was folded into its published
version when that was unambiguous (\cref{sec:same-work-guard,%
sec:version-dedup}).

\section{The check by meaning}
\label{sec:search-check}

For each question, its claim was put to OpenAlex's semantic search as a
plain question, and each of the 50 works returned was looked up in the
question's sessions (\texttt{scholar sessions check}).
A work that bore on the question and was in no session was taken as a
gap in the wordings, and a query by concept was added for it and run on
every database.
For the first question 29 of the 50 works were in no session; two
repair queries brought back 17 of them, and the 12 left outside were
grey literature, items from 2026 and one preprint on stopping.
For the second question 36 were missing; one repair query returned three
of the key comparisons, and the 33 left outside were mostly general
texts on systematic reviews.
For the third, 29 were missing; one repair query returned one of them, and
26 comparisons of bibliographic coverage stayed outside.
For the fourth, 4 were missing and no query was added; among them was a
comparison of screening by title with screening by title and abstract,
which one query's wording should have returned.
A further repair query for the third question was added because two of
the named studies had not come back on its concepts.

\section{Tiers and samples}
\label{sec:search-tiers}

Not every record was screened.
The first tier was:
\begin{enumerate}
\item the first 100 places of each query's OpenAlex and arXiv answers,
  fetched again with a limit of 100 into a companion session per
  question, because the version of Scholar used then sent Scopus and Web
  of Science no sort order and their answers came newest first, or
  mostly so (\cref{sec:measured-chapter});
\item the records returned by at least three databases, for the first
  question, and by at least five for the others;
\item the works of the check by meaning that were in the sessions;
\item the named studies.
\end{enumerate}
The rest of each question's records, pending in its session and not
decided in the companion session, was left unscreened, and 60 of them
were drawn at random (Python's \texttt{random.Random(20261123 + $q$)}
for question $q$, over the sorted paper ids) and screened as the first
tier was.
For the second and third questions, 60 records of the group returned by
three or more databases were also drawn (seed 20261123) and screened.
\Cref{tab:search-tiers} gives the counts; the intervals are the 95\,\%
modified Wilson intervals that Scholar computes
(\cref{sec:llm-sample-interval}).

\begin{table}
  \begin{sidecaption}{%
      The literature search in numbers: records returned (with
      duplicates across queries and databases), distinct records,
      records screened in the first tier, bearing records among the
      screened (supporting plus qualifying), cited works (read in full
      plus abstract only), and records left unscreened.
    }[tab:search-tiers]
    \begin{tabular}{@{}lrrrrrr@{}}
      \toprule
      & Returned & Distinct & First tier & Bearing & Cited & Unscreened \\
      \midrule
      \Cref{ch:stopping} & 2,264 & 931 & 454 & 23 + 30 & 8 + 3 & 477 \\
      \Cref{ch:abbreviated} & 8,627 & 4,123 & 679 & 12 + 26 & 4 + 11 & 3,384 \\
      \Cref{ch:databases} & 113,942 & 61,348 & 1,003 & 22 + 36 & 7 + 4 & 60,268 \\
      \Cref{ch:llm-screening} & 6,818 & 2,574 & 652 & 63 + 44 & 7 + 3 & 1,910 \\
      \bottomrule
    \end{tabular}
  \end{sidecaption}
\end{table}

\section{Screening and reading}
\label{sec:search-screening}

One AI agent screened every record, from the title and, for anything
about methods, the abstract; no classifier was used and no person checked
its verdicts or a sample of them.
Its decisions are recorded in the sessions as an agent's
(\texttt{scholar sessions decide -{}-source agent -{}-reasons-file}), each
with its own reason.
Each kept record has one category: it supports the claim, or it
qualifies or contradicts it.
Each discarded record is either off topic, with a criterion (not a
screening or search method, no abstract, or not a study), or adjacent (the
right field, another question), with the question's own label for what it
lacks (no evidence on stopping, no comparison, no evidence on retrieval,
no measurement of errors).
Since nobody reviewed the agent's verdicts, the tables of bearing records
in the following chapters mark a record as cited when the text cites it,
not because a person kept it.

The cited works were read in full where the publisher's copy, an open
copy or a preprint could be had; 30 of the 51 were read in full (six of
them as preprints of papers whose published versions were not obtained,
or that are preprints only), 21 only in their abstract, and each chapter
says which.
Twenty works were wanted and not obtained; each chapter names those that
would have borne on its question.

\section{Limits of the search}
\label{sec:search-limits}

One screener, an agent, decided everything; how often it erred is not
known.
Most of the evidence found is from health research, with two or three
studies per question from computing and two from education.
Several sources on language models are preprints.
Four Semantic Scholar answers of the third question are incomplete
(rate limits), and its Scopus answer to one query stopped at 4,991 of
18,141 records (HTTP 400); its OpenAlex answer to the same query holds
9,978 of the 26,809 records OpenAlex reported, and the session does not
mark it incomplete.
The first tier of the third question was larger than planned (1,003
records against 900).
Seven bearing records of the third question were decided in its companion
session, two of them records that the question's full session does not
hold; its table of bearing records therefore comes in two parts
(\cref{sec:databases-hits}).
